% BHU PYQ -- Home Science: Therapeutic Nutrition and Diet Counselling (2024-25 NEP)
\documentclass[11pt,a4paper]{article}
\usepackage[a4paper,top=2.5cm,bottom=2.5cm,left=2.8cm,right=2.8cm,headheight=14pt]{geometry}
\usepackage{amsmath,amssymb}
\usepackage[T1]{fontenc}
\usepackage[utf8]{inputenc}
\usepackage{lmodern,microtype,fancyhdr,enumitem,hyperref,lastpage,parskip}

\hypersetup{colorlinks=true,urlcolor=black,linkcolor=black,pdftitle={HSCSE11: Therapeutic Nutrition and Diet Counselling -- BHU NEP 2024-25}}

\pagestyle{fancy}\fancyhf{}
\renewcommand{\headrulewidth}{0.4pt}\renewcommand{\footrulewidth}{0.4pt}
\fancyhead[L]{\small \textit{Banaras Hindu University}}
\fancyhead[R]{\small \textit{HSCSE11: Therapeutic Nutrition (2024-25)}}
\fancyfoot[C]{\small Page \thepage\ of \pageref{LastPage}}

\newlist{parts}{enumerate}{2}
\setlist[parts,1]{label=(\alph*),leftmargin=*,itemsep=4pt,topsep=3pt}
\setlist[parts,2]{label=(\roman*),leftmargin=*,itemsep=2pt,topsep=2pt}

\newcommand{\pts}[1]{\hfill \textit{[#1]}}

\begin{document}

\begin{center}
  {\large\bfseries BANARAS HINDU UNIVERSITY}\\[2pt]
  {\normalsize B. A. / B. Sc. (Hons.) Semester I Examination, 2024-25 (NEP)}\\[6pt]
  \rule{0.8\linewidth}{0.5pt}\\[6pt]
  {\large\bfseries HOME SCIENCE}\\[4pt]
  {\large\bfseries Paper Code: HSCSE-11 : Therapeutic Nutrition and Diet Counselling}\\[4pt]
  {\normalsize Time: 2:30 Hours \hfill Full Marks: 70}\\[2pt]
  {\small REF NO. CE/Conf./D-I/110}
\end{center}
\rule{\linewidth}{0.4pt}
\smallskip

\noindent \textit{Instructions: Attempt all Sections A, B and C according to directions. Write your Roll No.\ at the top immediately on receipt of this question paper.}

\bigskip

\section*{SECTION -- A}
\noindent \textit{Note: Attempt any \textbf{three} questions from this section.}\pts{3 $\times$ 10 = 30}

\begin{enumerate}[label=\textbf{\arabic*.},leftmargin=*,itemsep=10pt]

  \item Define Therapeutic Nutrition. Discuss the principles and importance of therapeutic diet modifications.\pts{10}

  \item Explain the etiology, clinical symptoms, and dietary management of Diabetes Mellitus.\pts{10}

  \item Discuss dietary management in cardiovascular diseases with special reference to Hypertension and Atherosclerosis.\pts{10}

  \item Define Diet Counselling. Discuss the steps and communication skills required for effective diet counselling.\pts{10}

\end{enumerate}

\bigskip

\section*{SECTION -- B}
\noindent \textit{Note: Attempt any \textbf{four} questions from this section.}\pts{4 $\times$ 5 = 20}

\begin{enumerate}[label=\textbf{\arabic*.},leftmargin=*,start=5,itemsep=8pt]

  \item Differentiate between clear liquid diet, full liquid diet, and soft diet.\pts{5}

  \item Describe dietary modifications required for Peptic Ulcer patients.\pts{5}

  \item Explain the role of a dietitian in hospital and community settings.\pts{5}

  \item Write short notes on Enteral and Parenteral nutrition.\pts{5}

  \item Discuss dietary guidelines for Hepatitis patients.\pts{5}

\end{enumerate}

\bigskip

\section*{SECTION -- C}
\noindent \textit{Note: Short Answer / Objective questions (within 50 words each):}\pts{10 $\times$ 2 = 20}

\begin{enumerate}[label=\textbf{10.}(\roman*),leftmargin=*,itemsep=6pt]

  \item Define Glycemic Index (GI).\pts{2}
  \item What is DASH diet?\pts{2}
  \item Name two rich dietary sources of soluble fiber.\pts{2}
  \item Define Renal Failure.\pts{2}
  \item What is Sodium restriction in diet?\pts{2}
  \item Write the full form of RDA.\pts{2}
  \item Name the ketone bodies.\pts{2}
  \item Define Anorexia Nervosa.\pts{2}
  \item What is Oral Rehydration Therapy (ORT)?\pts{2}
  \item Define Steatorrhea.\pts{2}

\end{enumerate}

\end{document}
